% Default pandoc/md2pdf style: generic, unbranded baseline.
% Package docs (run `texdoc <name>`, e.g. `texdoc titlesec`):
%   fontspec-less fonts: texdoc tgheros / texdoc lmodern
%   titlesec (\titleformat, \titlespacing, sectionbreak hooks): texdoc titlesec
%   fvextra (breaklines/breakanywhere on Verbatim/code blocks): texdoc fvextra
%   xcolor (\definecolor, HTML model): texdoc xcolor
%   hyperref (\hypersetup options): texdoc hyperref
%   float (\floatplacement, the [H] specifier): texdoc float
% Pandoc's own manual (variables, --include-in-header, image attributes):
%   https://pandoc.org/MANUAL.html

\usepackage[T1]{fontenc}
\usepackage{tgheros}
\renewcommand{\familydefault}{\sfdefault}
\usepackage{lmodern}
\renewcommand{\ttdefault}{lmtt}
\usepackage{microtype}
\usepackage{parskip}
\usepackage{xcolor}
\usepackage{titlesec}
\usepackage{setspace}
\setstretch{1.15}

% Code blocks (```lang fences) don't wrap long lines by default and
% overflow the page. fvextra is what actually provides the
% breaklines/breakanywhere keys — plain fancyvrb doesn't have them.
\usepackage{fvextra}
\fvset{breaklines=true, breakanywhere=true, breaksymbolleft={}}

% Neutral, unbranded palette: dark slate for headings, mid gray for links,
% near-black for body text.
\definecolor{titlecolor}{HTML}{2D3E50}
\definecolor{subtitlecolor}{HTML}{3D5266}
\definecolor{linkcolor}{HTML}{595959}
\definecolor{bodycolor}{HTML}{262626}

\titleformat{\section}
  {\normalfont\LARGE\bfseries\color{titlecolor}}{\thesection}{1em}{}
\titleformat{\subsection}
  {\normalfont\large\bfseries\color{subtitlecolor}}{\thesubsection}{1em}{}

\AtBeginDocument{\color{bodycolor}}
\AtBeginDocument{\hypersetup{colorlinks=true, linkcolor=linkcolor, urlcolor=linkcolor}}

% Without this, images become LaTeX floats and can drift away from where
% they were written in the source (LaTeX moves floats to "wherever fits").
% Pin them exactly where they are instead.
\usepackage{float}
\floatplacement{figure}{H}
